\documentclass[12pt]{article}
\usepackage[utf8]{inputenc}
\usepackage[T1]{fontenc}
\usepackage[french]{babel}
\usepackage{amsmath}
\usepackage{chemfig}
\usepackage{float}
\usepackage{hyperref}

\title{Résumé : Structure Spatiale et Nomenclature des Molécules Organiques}
\author{}
\date{}

\begin{document}

\maketitle

\section{Introduction}
La structure spatiale des molécules organiques est essentielle pour comprendre leurs propriétés physiques et chimiques. Voici un résumé pour bien nommer et comprendre les molécules organiques simples.

\section{Électrons et Liaisons}
\subsection{Doublets Liants et Non Liants}
Les électrons dans une molécule peuvent former des doublets liants (entre deux atomes) ou des doublets non liants (sur un seul atome).

\subsection{Exemple : Méthanol}
\chemfig{[-2]H-[:0]C(-[:90]H)(-[:270]H)-[:180]O-[:270]H}

\section{Formes Géométriques des Molécules}
\subsection{Linéaire}
Les molécules linéaires ont un angle de 180° entre les atomes. Exemple : \chemfig{C(-[:90]H)=C(-[:270]H)-H}

\subsection{Trigonale Plane}
Les molécules trigonales planes ont un angle de 120°. Exemple : \chemfig{C(-[:90]H)(-[:210]H)=C(-[:330]H)-H}

\subsection{Tétraèdre}
Les molécules en forme de tétraèdre ont un angle de 109.5°. Exemple : \chemfig{C(-[:90]H)(-[:210]H)(-[:330]H)-H}

\section{Isomérie}
\subsection{Isomérie de Constitution}
Les isomères de constitution ont la même formule brute mais des formules semi-développées différentes.

\subsection{Isomérie de Position}
Les groupements fonctionnels sont identiques mais situés différemment sur la chaîne carbonée.

\subsection{Isomérie de Chaîne}
Les chaînes carbonées sont différentes.

\subsection{Isomérie de Fonction}
Les groupes fonctionnels sont différents.

\section{Stéréoisomérie}
\subsection{Énantiomères}
Les énantiomères sont des images l'un de l'autre dans un miroir et ne sont pas superposables.

\subsection{Diastéréoisomères}
Les diastéréoisomères ne sont pas des images l'un de l'autre dans un miroir.

\subsection{Isomères Z/E}
Les isomères Z/E sont déterminés par la position des substituants autour d'une double liaison.

\section{Nomenclature des Molécules Organiques}
\subsection{Alcanes}
Les alcanes sont nommés en fonction du nombre de carbones dans la chaîne principale.

\begin{center}
\begin{tabular}{|c|c|c|}
\hline
n & Formule brute & Nom de l'alcane linéaire \\ \hline
1 & $\mathrm{CH_4}$ & Méthane \\ \hline
2 & $\mathrm{C_2H_6}$ & Éthane \\ \hline
3 & $\mathrm{C_3H_8}$ & Propane \\ \hline
4 & $\mathrm{C_4H_{10}}$ & Butane \\ \hline
5 & $\mathrm{C_5H_{12}}$ & Pentane \\ \hline
\end{tabular}
\end{center}

\subsection{Groupes Caractéristiques}
Les groupes caractéristiques déterminent la famille de la molécule.

\begin{center}
\begin{tabular}{|c|c|c|}
\hline
GROUPE CARACTERISTIQUE & FAMILLE (Fonction) & FORMULE DE LA FAMILLE \\ \hline
Groupe hydroxyle —OH & ALCOOL —anol & R-OH \\ \hline
Groupe carbonyle —C— & ALDÉHYDE —al & R-C-H \\ \hline
 & CÉTONE —one & R1-C-R2 \\ \hline
Groupe carboxyle —C-OH & ACIDE CARBOXYLIQUE & R-C-OH \\ \hline
Groupe amino —NH2 & AMINE —amine & R-NH2 \\ \hline
\end{tabular}
\end{center}

\section{Configuration Absolue (CIP)}
\subsection{Ordre de Priorité}
Les substituants sont ordonnés par numéro atomique.

\subsection{Exemple}
Pour un carbone asymétrique, si la séquence des substituants 1 à 3 est dans le sens des aiguilles d'une montre, la configuration est R (sinon S).

\section{Conclusion}
La nomenclature et la structure spatiale des molécules organiques sont cruciales pour comprendre leurs propriétés et réactions. Utilisez ces règles pour nommer correctement les molécules organiques simples.

\end{document}
